\chapter{セットアップ手順まとめ（チートシート）}

\section{最初の1回だけ行う作業}

\begin{wslbox}[(1) uv のインストール]
\begin{lstlisting}[style=bash]
$ curl -LsSf https://astral.sh/uv/install.sh | sh
\end{lstlisting}
\end{wslbox}

ターミナルを閉じて開き直す（または \cmd{source \textasciitilde/.bashrc}）。

\begin{wslbox}[(2) 課題リポジトリの取得]
\begin{lstlisting}[style=bash]
$ cd ~
$ git clone \
  https://github.com/shiru2/NMR-lab-kadai.git
\end{lstlisting}
\end{wslbox}

\begin{wslbox}[(3) 仮想環境の作成・パッケージのインストール]
\begin{lstlisting}[style=bash]
$ cd \
  ~/NMR-lab-kadai/2026_B4_assignment
$ uv sync
\end{lstlisting}
\end{wslbox}

\section{毎回の作業}

\textbf{VS Code（推奨）}：VS Code で
\cmd{2026\_B4\_assignment} フォルダを開くと、
カーネルが自動で \cmd{.venv} に設定されます。

\begin{wslbox}[VS Code でフォルダを開く（ターミナルから）]
\begin{lstlisting}[style=bash]
$ code \
  ~/NMR-lab-kadai/2026_B4_assignment
\end{lstlisting}
\end{wslbox}

\textbf{JupyterLab（任意）}：

\begin{wslbox}[JupyterLab を使う場合]
\begin{lstlisting}[style=bash]
$ cd \
  ~/NMR-lab-kadai/2026_B4_assignment
$ uv run jupyter lab
\end{lstlisting}
\end{wslbox}

ブラウザで \cmd{http://localhost:8888} を開いて作業する。
終了するときはターミナルで \key{Ctrl} + \key{C}。

\section{よく使うコマンドまとめ}

\begin{center}
\begin{tabularx}{\linewidth}{lX}
  \toprule
  コマンド & 内容 \\
  \midrule
  \cmd{uv --version} & uv のバージョン確認 \\
  \cmd{uv sync} & 仮想環境を作成・更新する \\
  \cmd{uv run python スクリプト名} & Python スクリプトを実行する \\
  \cmd{uv run jupyter lab} & JupyterLab を起動する（任意） \\
  \cmd{source .venv/bin/activate} & 仮想環境を手動で有効化する \\
  \cmd{uv pip list} & インストール済みパッケージを確認する \\
  \cmd{uv add パッケージ名} & 新しいパッケージを追加する \\
  \cmd{git pull} & 課題ファイルを最新版に更新する \\
  \cmd{uv self update} & uv 自体をアップデートする \\
  \bottomrule
\end{tabularx}
\end{center}

\section{トラブル時の確認フロー}

課題ノートブックが動かないときは以下の順番で確認してください。

\begin{enumerate}
  \item \cmd{uv sync} を実行してパッケージが揃っているか確認する
  \item VS Code で \cmd{2026\_B4\_assignment} フォルダが
    開かれているか確認する（サブフォルダだけでは不十分）
  \item カーネルエラーが出たら \cmd{rm -rf .venv \&\& uv sync} を試す
  \item それでも解決しない場合は第~\ref{chap:errors}章を参照する
\end{enumerate}

\vfill

\begin{center}
\small
参考：uv 公式ドキュメント \url{https://docs.astral.sh/uv/}
\end{center}
